\documentclass{article}

\title{Summary of Synthetic Tabular Data Generation Surveys}
\author{Máté Baumann}
\date{March 2026}

\begin{document}

\maketitle

\section{Survey Summary}

Synthetic tabular data generation has become increasingly important because many machine learning tasks depend on structured data, while real tabular datasets are often scarce, expensive to collect, or difficult to share due to privacy regulations and access limitations. Both surveys describe synthetic data as a practical solution for data sharing, augmentation, and downstream analysis. They also show that synthetic data is useful not only for privacy-sensitive settings, but also for data augmentation, lower data collection costs, and limited data availability. \cite{chal,shi}.

At the same time, both papers emphasize that tabular data generation is a difficult problem. Tabular datasets often combine numerical and categorical features, contain missing values and imbalanced classes, and include complex dependencies between columns. Because of this, a successful model must preserve not only individual feature distributions, but also meaningful relations across the table. This makes tabular synthesis more challenging than it may first appear \cite{chal,shi}.

To address these issues, the surveys discuss three broad method families: traditional methods, diffusion-based methods, and LLM-based methods. Traditional approaches include GAN-based models such as CTGAN, which are especially noted for handling imbalanced categorical variables in tabular data. Diffusion models such as TabDDPM have gained attention because they often produce high-quality synthetic samples, while newer LLM-based methods such as GReaT extend the field by modeling tabular rows through language-model techniques. Together, these families show how the field has moved from earlier deep generative models toward more recent diffusion and language-model-based approaches \cite{chal,shi}.

Both surveys also identify privacy as one of the central open problems in synthetic tabular data generation. They stress that realistic or useful synthetic data is not automatically private, and that evaluation must consider privacy risks alongside fidelity and machine learning utility. This leads to a continuing privacy--utility trade-off, where stronger privacy protection can reduce usefulness, while highly realistic generation can increase the risk of information leakage. Overall, the two surveys present a fast-developing field with strong recent progress, but also with clear remaining challenges in privacy, evaluation, and method selection \cite{chal,shi}.


\bibliographystyle{plain}
\bibliography{references}

\end{document}
